\section{Robustness and Extensions}
\label{sec:robustness}

We subject the conditional-association baselines and the
institutional-identification estimates to five groups of checks:
FL-definition variants (IQR multiplier, win-rate cutoff, temporal
window); sample and specification variants (unrestricted sample,
homogeneous subsamples, item$\times$year FE, matching);
alternative clustering; sensitivity to unobservables
(\citealt{cinelli2020making, oster2019unobservable}); and
oversight heterogeneity. We also report two designs whose
identifying assumptions fail in this setting and explain why we
treat them as transparency artifacts rather than as estimates.
Table~\ref{tab:robustness_summary} collects the point
estimates with the assumption each estimator imposes.

\begin{table}[htbp]
\centering
\caption{Robustness summary: FL price coefficient across checks}
\label{tab:robustness_summary}
\small
\begin{tabular}{lcccll}
\hline\hline
Check & Coeff.\ & SE & $N$ & Assumption & Note \\
\hline
\multicolumn{6}{l}{\textit{Baseline and matching}} \\
OLS (item + year + PBU FE) & 0.064 & 0.020 & 1{,}654{,}401 & Sel.\ on obs.\ + FE & Baseline \\
CEM & 0.077 & 0.024 & 969{,}751 & Sel.\ on obs. & \S\ref{sec:results_ols} \\
IPW & 0.055 & 0.021 & 830{,}194 & Sel.\ on obs. & \S\ref{sec:results_ols} \\
Cross-fit average & 0.036 & --- & --- & Sel.\ on obs.\ + FE; no mech.\ link & Out-of-sample \\
\hline
\multicolumn{6}{l}{\textit{FL definition variants}} \\
IQR 1.0$\times$ & 0.079 & 0.023 & 1{,}654{,}401 & Sel.\ on obs.\ + FE & 3{,}442 FL firms \\
IQR 1.5$\times$ & \multicolumn{5}{l}{\textit{(= OLS baseline above)}} \\
IQR 2.0$\times$ & 0.060 & 0.022 & 1{,}654{,}401 & Sel.\ on obs.\ + FE & 2{,}093 FL firms \\
IQR 3.0$\times$ & 0.050 & 0.024 & 1{,}654{,}401 & Sel.\ on obs.\ + FE & 1{,}456 FL firms \\
\hline
\multicolumn{6}{l}{\textit{Specification variants}} \\
Item$\times$year FE & 0.074 & 0.022 & 1{,}654{,}401 & Sel.\ on obs.\ + FE & Tighter controls \\
Two-way clustering & 0.064 & 0.024 & 1{,}654{,}401 & Sel.\ on obs.\ + FE & Item + PBU \\
\hline
\multicolumn{6}{l}{\textit{Alternative treatments and diagnostics}} \\
Continuous (log max AL) & 0.022 & 0.005 & 1{,}654{,}401 & Sel.\ on obs.\ + FE & Per log-point \\
Decomposition (odd-yr FL, full sample) & 0.043 & 0.019 & 1{,}654{,}401 & Sel.\ on obs.\ + FE & Intermediate \\
Horse-race (FL + Imhof CV) & 0.084 & 0.023 & 1{,}654{,}401 & Sel.\ on obs.\ + FE & Suppression \\
IV (leave-one-out) & 0.194 & 0.077 & 1{,}654{,}401 & Excl.\ restriction & Measurement error \\
\hline\hline
\end{tabular}
\begin{tablenotes}
\small
\item \textit{Notes:} ``Sel.\ on obs.\ + FE'' = selection on observables absorbed by item, year, and PBU fixed effects. Cross-fit further removes any mechanical link between the FL rule and the price outcome. The IV is a measurement-error diagnostic; exclusion-restriction concerns (\S\ref{sec:results_ols}) keep it off the headline range. Staggered DiD (Callaway--Sant'Anna and stacked Cengiz et al.) and Oster $\hat{\delta}$ are excluded from this summary because their identifying assumptions fail in this setting; they are reported transparently in \S\ref{sec:robustness}, \ref{app:did}, and \ref{app:robustness}.
\end{tablenotes}
\end{table}


\paragraph{Threshold sensitivity.}
The FL coefficient declines monotonically with stricter
thresholds: $0.079$ (IQR $1.0\times$), $0.064$ (baseline
$1.5\times$), $0.060$ ($2.0\times$), $0.050$ ($3.0\times$).
Two readings are consistent with the pattern. \emph{(i)}
Classification noise: stricter thresholds shrink the treated
group, leaving fewer FL-present tenders and more misclassified
controls. \emph{(ii)} Marker bite: the marginal FL firms
(10--14 tenders) sit in markets with larger cartel markups
while tail firms face more detection pressure. The noise reading
is the simpler explanation and we adopt it. A two-dimensional
sensitivity analysis varying the IQR multiplier (1.0--3.0$\times$)
and the win-rate cutoff (0--5\%) yields significant
coefficients across all 36 cells
(Figure~\ref{fig:threshold_heatmap}, \ref{app:robustness});
the result does not hinge on the FL definition.

\paragraph{Sensitivity to unobservables.}
The \citet{cinelli2020making} sensitivity analysis yields
$RV_{q=1} = 17.5\%$: an unobserved confounder would need to
explain at least $17.5\%$ of residual variation in
\emph{both} FL presence and log prices to nullify the
coefficient---stringent given the item, year, and PBU
fixed-effects structure. Sensitivity contour plots appear in
\ref{app:robustness}. The \citet{oster2019unobservable}
bound is uninformative here because adding PBU fixed effects
moves $R^2$ from $0.879$ to $0.886$; the denominator of the
Oster formula is near zero and $\hat{\delta} = 261.6$ is
mechanically large. The Cinelli--Hazlett bound, which does not
depend on the $R^2$ gap, is the more informative robustness
statement for this design.

\paragraph{Oversight heterogeneity.}
The FL price coefficient varies steeply with PBU size: $0.214$
in the smallest quartile, $0.098$ in Q2, $0.045$ in Q3, and
$0.017$ in Q4 (Table~\ref{tab:regime_oversight},
\ref{app:robustness}). The 12.5$\times$ gradient between Q1
and Q4 is the single largest source of heterogeneity in the
data---larger than the preg\~ao--convite gap---and runs in the
direction the cover-bidding framework implies
($\partial m^*/\partial \theta_k < 0$, \ref{app:proofs}):
smaller PBUs with less oversight capacity see more FL activity
and higher associated price gaps. The alternative reading is
that smaller PBUs have thinner item--PBU cells, so item FE
absorb less within-market variation and the coefficient picks
up residual heterogeneity. The monotonic gradient across all
four quartiles is hard to generate from pure FE noise, but we
cannot fully rule it out.

\paragraph{Designs we attempted and rejected.}
Two estimators that the recent staggered-DiD literature
recommends fail their identifying assumptions in this setting,
and we report them only to document the failures.
\citet{callaway2021difference} doubly-robust staggered DiD
yields $\mathrm{ATT} = 0.014$ (SE $0.039$); the
\citet{cengiz2019effect} stacked design yields
$\mathrm{ATT} = -0.006$ (SE $0.014$;
Table~\ref{tab:stacked_did}, \ref{app:did}). The stacked 95\%
CI $(-0.034, 0.022)$ excludes the OLS baseline of $0.064$. The
discrepancy is mechanical: positive pre-event coefficients at
$t = -2$ and $t = -3$ (Figure~\ref{fig:event_study},
\ref{app:did}) indicate that FL entry is endogenous to market
conditions, violating the parallel-trends assumption both
estimators require. We declare the designs \emph{invalid} for
this setting, not merely underpowered, and exclude them from
Table~\ref{tab:robustness_summary}. \citet{rambachan2023more}
sensitivity bounds confirm that the maximum pre-trend slope
under which the conclusion would hold is implausibly small.
\citet{oster2019unobservable} $\hat{\delta} = 261.6$ is
similarly degenerate (above). The right reading of these
exercises is that staggered-DiD identification requires
exogenous treatment timing the FL setting does not provide,
and the conditional-association range we report is robust to
this constraint, not despite it.

\paragraph{Illustrative welfare.}
\ref{app:extensions} reports back-of-the-envelope welfare
calculations that mechanically apply the FL price coefficients
to FL-present spending: roughly $0.3$--$0.9\%$ of total BEC
spending depending on which point estimate is used. The
calculations \emph{assume} causality the design does not
establish; they should be read as rough benchmarks, not policy
estimates. A causal welfare quantification awaits replication
in a setting with stronger institutional identification (e.g.,
post-2021 BEC under Lei 14.133, or external replication in
ComprasNet or OECD platforms; see Section~\ref{sec:conclusion}).
